\documentclass[11pt]{article}
\usepackage[a4paper,margin=23mm]{geometry}
\usepackage{amsmath,amssymb,amsthm,mathtools,booktabs,array}
\usepackage{fontspec,microtype,xurl}
\usepackage[colorlinks=true,linkcolor=blue,urlcolor=blue]{hyperref}
\newcommand{\R}{\mathbb R}
\newcommand{\PP}{\mathcal P}
\newcommand{\norm}[1]{\left\lVert#1\right\rVert}
\newcommand{\diag}{\operatorname{diag}}
\newcommand{\rank}{\operatorname{rank}}
\newcommand{\Var}{\operatorname{Var}}
\newcommand{\code}[1]{\texttt{\detokenize{#1}}}
\newtheorem{theorem}{Theorem}[section]
\newtheorem{lemma}[theorem]{Lemma}
\newtheorem{proposition}[theorem]{Proposition}
\newtheorem{corollary}[theorem]{Corollary}
\theoremstyle{remark}\newtheorem{remark}[theorem]{Remark}
\setlength{\emergencystretch}{3em}
\title{Four focused re-derivations in polynomial approximation}
\author{Mathematical handoff for NLAlib}
\date{13 October 2026}
\begin{document}
\maketitle
\noindent This document contains only rational equioscillation, sharp logarithmic
Lebesgue estimates (including the universal interpolation lower bound), and
bounded-variation coefficient and quadrature estimates. All functions and
polynomials are real unless stated otherwise. Write $\PP_d$ for the real
polynomials of degree at most $d$, including zero, and use the uniform norm
on $[-1,1]$.

\paragraph{What has and has not been simplified.}
The rational proof separates compactness, cancellation, degree defects, and
a finite signed annihilator. The logarithmic estimates use explicit finite
cotangent sums and ordinary scalar integrals. The variation proof uses step
functions and a finite coefficient recurrence. Quadrature uses exact folded
moments and gives an explicit eventual threshold smaller than the earlier
$14\nu4^\nu$ threshold.
For arbitrary interpolation nodes, a purely finite replacement for the
global comparison theorem has \emph{not} been established here. Instead,
Section~\ref{sec:universal} supplies a complete shorter continuation argument:
compactness and finitely many local inverse charts replace the general
proper-map/global-inverse theorem. It is still a global argument. No Lean
formalization or tested Lean implementation is claimed.
\tableofcontents

\section{Rational equioscillation with exceptional cases separated}
\label{sec:rational}
Fix integers $m,n\ge0$. Let $\mathcal R_{m,n}$ consist of rational functions
representable by $p/q$, with $p\in\PP_m$, $0\ne q\in\PP_n$, after removal of
removable singularities, and with no pole on $[-1,1]$. Every such function has
a reduced representative with $q>0$ on this interval. Uniqueness below is
uniqueness of the function, not of its numerator and denominator.

\subsection{Compactness and cancellation, before alternation}
\begin{lemma}[A fixed evaluation map controls coefficients]\label{lem:vander}
For $d+1$ fixed distinct points $z_0,\ldots,z_d$ in $[-1,1]$, a uniform bound
on $|p(z_i)|$ for $p\in\PP_d$ bounds the coefficient vector of $p$.
Coefficient convergence in $\PP_d$ implies uniform convergence on $[-1,1]$.
\end{lemma}
\begin{proof}
The evaluation matrix is the Vandermonde matrix $(z_i^j)_{i,j=0}^d$, whose
determinant is $\prod_{i<j}(z_j-z_i)\ne0$. Apply its fixed inverse. For the
second assertion, $\norm{\sum_{j=0}^d c_jx^j}_\infty\le\sum_{j=0}^d|c_j|$.
\end{proof}

\begin{lemma}[Local boundedness forces cancellation]\label{lem:cancel}
Suppose $q\ne0$ and $|p(x)|\le C|q(x)|$ throughout $[-1,1]$. If $p\ne0$,
every zero of $q$ in the interval has multiplicity at most its multiplicity
as a zero of $p$. Consequently canceling the polynomial gcd gives a fraction
whose denominator has no zero on the interval. Degrees do not increase.
If $p=0$, use the reduced representative $0/1$ instead.
\end{lemma}
\begin{proof}
At a zero $z$ factor $p=(X-z)^u p_0$ and $q=(X-z)^v q_0$, with
$p_0(z)q_0(z)\ne0$. If $u<v$, then $|p(x)/q(x)|$ tends to infinity as
$x\to z$ through interval points different from $z$, contrary to its bound.
This uses a one-sided limit at $z=\pm1$. Polynomial gcd cancellation removes
all these denominator zeros. The remaining denominator is continuous and
nonzero, hence has a constant sign; reverse both signs if necessary.
\end{proof}

\begin{theorem}[Existence]\label{thm:rexist}
Every $f\in C[-1,1]$ has a best approximant in $\mathcal R_{m,n}$.
\end{theorem}
\begin{proof}
Let $E=\inf_{r\in\mathcal R_{m,n}}\norm{f-r}_\infty$. It is finite and
$0\le E\le\norm f_\infty$, since zero is admissible. Choose a minimizing
sequence $p_j/q_j$, with $q_j>0$, errors at most $\norm f_\infty+1$, and
normalize the Euclidean norm of the coefficient vector of $q_j$ to one.
Compactness gives a subsequence converging coefficientwise to $q$ with the
same coefficient norm. Thus $q\ne0$ and $q\ge0$ on the interval.

The normalization gives $\norm{q_j}_\infty\le\sqrt{n+1}$ and
$\norm{p_j/q_j}_\infty\le2\norm f_\infty+1$. Hence the $p_j$ are uniformly
bounded. Lemma~\ref{lem:vander} and another subsequence give $p_j\to p$
coefficientwise. Passing to the limit gives
\begin{equation}\label{eq:limit}
 |f(x)q(x)-p(x)|\le E q(x),\qquad
 |p(x)|\le(\norm f_\infty+E)|q(x)|.
\end{equation}
If $p\ne0$, apply Lemma~\ref{lem:cancel}; if $p=0$, take $0/1$.
Off the finitely many zeros of $q$, the resulting continuous rational
function has error at most $E$ by \eqref{eq:limit}. These points are dense,
so the same inequality holds at the canceled zeros by continuity.
\end{proof}

\subsection{Degree defects are exactly a finite rank calculation}
For nonzero reduced $r=p/q$, put $\mu=\deg p$, $\eta=\deg q$, and
\begin{equation}\label{eq:defect}
 d=\min(m-\mu,n-\eta),\qquad
 D=m+n-d=\max(m+\eta,n+\mu).
\end{equation}
For $r=0$, use $p=0,q=1$ and \emph{define $D=m$}. In particular, do not
substitute degree zero for the zero polynomial in \eqref{eq:defect}.

\begin{lemma}[Surjectivity of the cross-numerator map]\label{lem:rank}
For a nonzero reduced fraction, the map
\[
 \mathcal L:\PP_m\times\PP_n\longrightarrow\PP_D,
 \qquad (P,Q)\longmapsto Pq-pQ
\]
is onto. Its kernel is precisely $\{(ph,qh):h\in\PP_d\}$.
For the zero fraction the needed onto map is simply $P\mapsto P$ on $\PP_m$.
\end{lemma}
\begin{proof}
Coprimality and $Pq=pQ$ give $P=ph,Q=qh$ with the same polynomial $h$.
The degree constraints give exactly $h\in\PP_d$. This parameterization is
injective since $q\ne0$. Rank--nullity gives image dimension
$(m+1)+(n+1)-(d+1)=D+1$, the dimension of $\PP_D$.
\end{proof}

\subsection{Two scalar lemmas for necessity}
\begin{lemma}[Separate the extremal sign blocks]\label{lem:signs}
Let $e\in C[-1,1]$, $E=\norm e_\infty>0$. If the sets
$K_\pm=\{x:e(x)=\pm E\}$ contain no increasing alternating list of length
$D+2$, there is $h\in\PP_D$ with $e(x)h(x)>0$ on $K_+\cup K_-$.
\end{lemma}
\begin{proof}
The two compact sign sets are disjoint, so their distance is positive when
both are nonempty. Starting at their smallest point, repeatedly select the
first later point of the opposite sign. Compactness gives each first point;
positive separation makes the process finite. The hypothesis limits it to
$D+1$ points. Between two successive selected points, the extremal points
of the first sign have a supremum strictly below the next selected point.
Choose a separator in this gap. The product of $X-z$ over all separators,
with its overall sign chosen suitably, has the required signs. Its degree
is the number of separators, at most $D$. A single sign block uses a
constant polynomial.
\end{proof}

\begin{lemma}[A uniform signed descent]\label{lem:descent}
Let $v_t\to v$ uniformly as $t\downarrow0$, with continuous functions.
If $e(x)v(x)>0$ whenever $|e(x)|=E=\norm e_\infty>0$, then
$\norm{e-tv_t}_\infty<E$ for all sufficiently small $t>0$.
\end{lemma}
\begin{proof}
On the compact extremal set, $|v|$ has positive minimum $c$. Choose a
neighborhood $U$ of this set on which $e,v$ have the same sign,
$|v|\ge3c/4$, and $|e|>E/2$. Require $\norm{v_t-v}_\infty<c/4$ and a
uniform bound $\norm{v_t}_\infty\le B$. On the compact complement there
is a margin $|e|\le E-\varepsilon$; omit this condition if the complement
is empty. For $tB<\min(E/2,\varepsilon/2)$ there is no sign overshoot on
$U$, where $|e-tv_t|\le E-tc/2$; outside $U$, the bound is
$E-\varepsilon/2$.
\end{proof}

\subsection{One finite annihilator proves optimality and uniqueness}
For $x_0<\cdots<x_{D+1}$ set
\[
 w_i=\frac1{\prod_{j\ne i}(x_i-x_j)},\qquad
 b_i=\frac{|w_i|}{\sum_j|w_j|}>0.
\]
Leading-coefficient comparison in the Lagrange formula gives
\begin{equation}\label{eq:annihilator}
 \sum_i w_i H(x_i)=0\quad(H\in\PP_D),\qquad
 \sum_i b_i\sigma(-1)^i H(x_i)=0\quad(\sigma=\pm1).
\end{equation}
Indeed $\operatorname{sgn}w_i=(-1)^{D+1-i}$.

\begin{theorem}[Rational equioscillation, including zero]\label{thm:rational}
For $r\in\mathcal R_{m,n}$ define $D$ as above. Then $r$ is best for $f$
if and only if there are $D+2$ ordered points and $\sigma=\pm1$ such that
\[
 (f-r)(x_i)=\sigma(-1)^i\norm{f-r}_\infty.
\]
The best function is unique. More generally, $D+2$ alternating errors
of magnitudes at least $\varepsilon>0$ certify that every competitor has
uniform error at least $\varepsilon$.
\end{theorem}
\begin{proof}
For necessity at positive error, a failure of alternation supplies $h$
from Lemma~\ref{lem:signs}. For $r\ne0$, Lemma~\ref{lem:rank} supplies
$P,Q$ with $Pq-pQ=h$. For small $t>0$, $q+tQ$ stays strictly positive,
and
\[
 r_t-r=\frac{p+tP}{q+tQ}-\frac pq
       =t\frac{h}{q(q+tQ)}.
\]
The quotient after $t$ converges uniformly to $h/q^2$, so
Lemma~\ref{lem:descent} contradicts optimality. For $r=0$, use the
admissible polynomial perturbation $r_t=th$. Zero error permits any
$D+2$ distinct interval points.

For sufficiency compare $r=p/q$ with $s=P/Q$, where $q,Q>0$. Its
cross-numerator $H=Pq-pQ$ has degree at most $D$, even when the competitor
has a different defect; for the zero baseline it is $P\in\PP_m$.
An error strictly less than $\varepsilon$ forces
$\sigma(-1)^i(s-r)(x_i)>0$. Positive denominators transfer the signs to
$H$, contradicting \eqref{eq:annihilator}.

For uniqueness let $r$ be best with positive error $E$ and use its
alternating list. Another best $s$ satisfies
$\sigma(-1)^i H(x_i)\ge0$. Equation~\eqref{eq:annihilator} is a sum
of nonnegative quantities with strictly positive weights; hence every
$H(x_i)=0$. A polynomial of degree at most $D$ with $D+2$ distinct roots
is zero. Therefore $s=r$. At error zero both functions already equal $f$.
Together with Theorem~\ref{thm:rexist}, this proves existence and uniqueness.
\end{proof}

\begin{remark}
This uniqueness proof needs neither a midpoint of rational functions nor
an additional defect argument for that midpoint. For $r=0$, the correct
alternation count is always $m+2$, independently of the denominator budget
$n$. All degree budgets $m=0$ or $n=0$ are included.
\end{remark}

\section{Sharp logarithmic bounds from finite kernels}
\label{sec:kernels}
Write $T_k(\cos\theta)=\cos(k\theta)$. The coefficient convention used
throughout the document is
\begin{equation}\label{eq:coeff}
 a_0(f)=\frac1\pi\int_0^\pi f(\cos\theta)\,d\theta,\qquad
 a_k(f)=\frac2\pi\int_0^\pi f(\cos\theta)\cos(k\theta)\,d\theta\quad(k\ge1).
\end{equation}
Let $I_n$ interpolate at $\cos(j\pi/n)$, $0\le j\le n$, for $n\ge1$,
and let $S_nf=\sum_{k=0}^n a_k(f)T_k$.

\subsection{The finite cotangent inequality for Lobatto interpolation}
\begin{theorem}\label{thm:lobatto}
For $n\ge1$,
\[
 \norm{I_n}=\Lambda_n\le
 C_n:=\frac1n\sum_{j=0}^{n-1}\cot\frac{(j+1/2)\pi}{2n}
 \le1+\frac2\pi\log(n+1).
\]
\end{theorem}
\begin{proof}
Define the finite trigonometric cardinal kernel
\[
 K_n(t)=\frac1{2n}\left(1+2\sum_{k=1}^{n-1}\cos kt+\cos nt\right)
       =\frac{\sin(nt)\cot(t/2)}{2n}.
\]
The trigonometric polynomial defines the removable values. Its translates
at $t_j=j\pi/n$, $0\le j<2n$, satisfy $K_n(t_i-t_j)=\delta_{ij}$.
For symmetric data, their interpolation sum is even, hence a polynomial
of degree at most $n$ in $\cos t$, and equals $I_n$. Thus
$\Lambda_n\le\max_t\sum_j|K_n(t-t_j)|$.

By translation and reflection take $t=\alpha\pi/n$, $0<\alpha<1$.
With $h=\pi/(2n)$ the last sum is
\begin{equation}\label{eq:cotpair}
 \frac{\sin\pi\alpha}{2n}\sum_{j=0}^{n-1}
 \{\cot(h(j+\alpha))+\cot(h(j+1-\alpha))\}.
\end{equation}
Each summand, including its sine multiplier, is bounded by
$2\cot(h(j+1/2))$. To prove this set $s=|\alpha-1/2|$,
$u=h(j+1/2)$, $v=hs$. The addition formula
\[
 \cot(u+v)+\cot(u-v)=\frac{\sin2u}{\sin^2u-\sin^2v}
\]
reduces the assertion to
$\sin^2(hs)\le2\sin^2(\pi s/2)\sin^2u$.
For $s>0$, the ratio $\sin(hs)/\sin(\pi s/2)$ is increasing in $s$:
its logarithmic derivative is nonnegative because $z\cot z$ is decreasing
on $(0,\pi/2]$. The latter follows from $\sin z\cos z\le z$.
At $s=1/2$ the ratio is $\sqrt2\sin(h/2)\le\sqrt2\sin u$, proving
the claim. At $s=0$ it is immediate. Equality in \eqref{eq:cotpair}
holds at $\alpha=1/2$, so its maximum is $C_n$. At grid points it is one,
which is no larger: the cardinal functions sum to one.

Convexity of $\cot$ on $(0,\pi/2]$, applied on the cells
$[jh,(j+1)h]$, $1\le j<n$, gives
\[
 C_n\le\frac1n\cot(h/2)+\frac1{nh}\int_h^{\pi/2}\cot u\,du
 \le\frac4\pi+\frac2\pi\log\frac1{\sin h}.
\]
Furthermore $\sin h\ge h-h^3/6\ge\pi/[2(n+1)]$; the last inequality
is $\pi^2(n+1)\le24n^2$. Hence
\[
 C_n\le\frac2\pi\log(n+1)+\frac2\pi(2+\log(2/\pi))
       \le\frac2\pi\log(n+1)+1.
\]
For the numerical inequality, $\log x\ge2(x-1)/(x+1)$ for $x\ge1$
by differentiation. At $x=\pi/2$ it suffices that
$\pi^2+2\pi-16>0$, which follows from $\pi>25/8$.
The case $n=1$ also has $C_1=1$.
\end{proof}

\subsection{The finite Dirichlet kernel and the sharp logarithmic coefficient}
\begin{theorem}\label{thm:projection}
Let $H_n=\sum_{j=1}^n1/j$. For $n\ge1$,
\[
 \frac4{\pi^2}H_n\le\norm{S_n}
 =\frac1{2\pi}\int_{-\pi}^{\pi}|D_n(t)|\,dt
 \le3+\frac4{\pi^2}\log(n+1),
 \qquad D_n(t)=1+2\sum_{k=1}^n\cos kt.
\]
For $n=0$, $\norm{S_0}=1$.
\end{theorem}
\begin{proof}
The finite product-to-sum formula gives, for $g(t)=f(\cos t)$,
\[
 (S_nf)(\cos\theta)=\frac1{2\pi}\int_{-\pi}^\pi D_n(\theta-t)g(t)\,dt.
\]
This proves the norm upper estimate by the integral. For the reverse
inequality put $d_n(x)=1+2\sum_{k=1}^nT_k(x)$ and
$f_\varepsilon=d_n/\sqrt{d_n^2+\varepsilon^2}$. This is a continuous
function bounded by one. Its image at $x=1$ differs from the kernel
integral by at most $\varepsilon$, because
$0\le |u|-u^2/\sqrt{u^2+\varepsilon^2}\le\varepsilon$.

The finite geometric identity gives
$D_n(t)=\sin((n+1/2)t)/\sin(t/2)$ away from its removable points.
For $0<t\le\pi$,
\[
 \frac2t\le\frac1{\sin(t/2)}\le\frac2t+\frac t4.
\]
For the upper bound use $\sin u\ge u-u^3/6\ge u/(1+u^2/2)$ on
$0\le u\le\pi/2<2$; multiplying the second inequality leaves
$u^3(4-u^2)/12\ge0$. Write $M=n+1/2$. Then
\[
 \norm{S_n}\le\frac2\pi\int_0^{M\pi}\frac{|\sin u|}{u}\,du+\frac\pi8,
 \qquad
 \int_0^{M\pi}\frac{|\sin u|}{u}\,du\le\pi+\frac2\pi H_n.
\]
The first full period is bounded by $\pi$ and each later full period
$[j\pi,(j+1)\pi]$ by $2/(j\pi)$; extending the final half period is
harmless. Now $H_n\le1+\log(n+1)$ and
$\pi/8+4/\pi^2<11/28+4/9<1$.

For the lower bound retain the first $n$ complete periods and use
$1/u\ge1/((j+1)\pi)$ on period $j$. Their sine integrals are two, so
$\norm{S_n}\ge(4/\pi^2)H_n$. This also proves sharpness of the
logarithmic coefficient, rather than only an upper estimate.
\end{proof}

\begin{corollary}[Operator consequences]
With $E_n(f)=\inf_{p\in\PP_n}\norm{f-p}_\infty$,
\[
 \norm{f-I_nf}_\infty\le(2+(2/\pi)\log(n+1))E_n(f),\qquad
 \norm{f-S_nf}_\infty\le(4+(4/\pi^2)\log(n+1))E_n(f).
\]
\end{corollary}
\begin{proof}
Both operators fix $\PP_n$. For any such operator $A$ and polynomial $p$,
$f-Af=(f-p)-A(f-p)$; take the infimum after the triangle inequality.
\end{proof}

\section{The exact universal lower bound: finite certificates and a smaller global lemma}
\label{sec:universal}
For distinct nodes $X=(x_0,\ldots,x_n)$ in $[-1,1]$, let
\[
 \ell_j^X(x)=\prod_{r\ne j}\frac{x-x_r}{x_j-x_r},\quad
 \mathcal L_X(x)=\sum_j|\ell_j^X(x)|,\quad
 \Lambda_n(X)=\max_{x\in[-1,1]}\mathcal L_X(x).
\]
The interpolation operator has norm $\Lambda_n(X)$: the upper bound is the
triangle inequality, and the reverse bound follows by assigning at the
nodes the signs of $\ell_j^X(x_*)$ at a maximizing point, then extending
those data piecewise affinely with uniform norm one.
We prove
\[
 \Lambda_n(X)\ge\frac2\pi\bigl(\log(n+1)+\gamma+\log(4/\pi)\bigr).
\]
The finite estimates below preserve the additive constant as well as
$2/\pi$. The passage from one grid to every grid still uses compactness
and local inverses. The last subsection proves exactly that passage;
it neither assumes a general global inverse theorem nor claims that
local Jacobian signs alone imply the result.

Write $N=n+1$ for the number of nodes and put
\[
 B_N=\frac2\pi\left(\log N+\gamma+\log\frac4\pi\right),
 \qquad H_m=\sum_{j=1}^{m}\frac1j,
\]
where $\gamma=\lim_m(H_m-\log m)$.
For the $N$ Chebyshev roots
$\xi_j=\cos((j+\tfrac12)\pi/N)$, $0\le j<N$, let
$\mathcal L_N(x)$ denote their Lebesgue function.  Its order of indexing
does not affect the function.

\subsection{Evaluate at extrema instead of finding the local maxima}
\label{polyretry:cotangent-samples}

For $N\ge2$ and $1\le k<N$ set $y_k=\cos(k\pi/N)$ and
$S_k=\mathcal L_N(y_k)$.  There is exactly one such point strictly
between each pair of consecutive roots.  The cardinal formula obtained
by differentiating $T_N$ gives
\begin{equation}\label{polyretry:sample}
 S_k=\frac1N\sum_{j=0}^{N-1}
 \frac{\sin((j+\tfrac12)\pi/N)}
 {|\cos(k\pi/N)-\cos((j+\tfrac12)\pi/N)|}.
\end{equation}
Indeed $|T_N(y_k)|=1$ and
$|T_N'(\xi_j)|=N/\sin((j+\tfrac12)\pi/N)$.

Here is a finite calculation locating the smallest \emph{sample},
which is enough for a lower bound on all interval maxima.  Put
$h=\pi/(2N)$ and
$C_r=\sum_{j=0}^{r-1}\cot((j+\tfrac12)h)$, $0\le r\le2N$.
The identity
\[
 \frac{\sin\theta}{\cos\phi-\cos\theta}
 =\frac12\left(\cot\frac{\theta-\phi}{2}
                  +\cot\frac{\theta+\phi}{2}\right)
\]
and $C_{N+r}=C_{N-r}$ give, by splitting the sum at $j=k$,
\begin{equation}\label{polyretry:sample-C}
 S_k=\frac{C_{N-k}+C_k-C_{2k}}N.
\end{equation}
Consequently
\[
 N(S_{k+1}-S_k)=2\cot((2k+1)h)
 -\cot((2k+\tfrac12)h)-\cot((2k+\tfrac32)h).
\]
For $1\le k<\lfloor N/2\rfloor$ this is nonpositive, since
$\cot''x=2\cot x\csc^2x>0$ on $(0,\pi/2)$.
Also $S_k=S_{N-k}$ by reflection.  Thus the smallest sample is
$S_m$, where $m=\lfloor N/2\rfloor$.
Using $\cot u-\tan u=2\cot(2u)$ in
\eqref{polyretry:sample-C} gives the explicit formula
\begin{equation}\label{polyretry:central-cot}
 S_m=\frac2N\sum_{j=0}^{m-1}
       \cot\frac{(j+\tfrac12)\pi}{N}
 +\begin{cases}
 0,&N=2m,\\[2pt]
 N^{-1}\tan\dfrac{\pi}{4N},&N=2m+1.
 \end{cases}
\end{equation}
In particular no theorem about the locations or monotonicity of the
\emph{true} local maxima is needed.

\subsection{The exact logarithmic constant by midpoint concavity}

Define $g(x)=\cot x-1/x$ on $(0,\pi/2]$, with $g(0)=0$.
This function is concave.  To check the only nontrivial sign,
\[
 g''(x)=\frac{2\cos x}{\sin^3x}-\frac2{x^3}\le0
 \quad\Longleftrightarrow\quad
 \sin^3x\ge x^3\cos x.
\]
For $0<x<\pi/2$, differentiate
$3\log(\sin x/x)-\log\cos x$.  The sign of its derivative is
the sign of
\[
 K(x)=x(2+\cos2x)-\tfrac32\sin2x.
\]
But $K(0)=0$ and
$K'(x)=4\sin x(\sin x-x\cos x)\ge0$, because
$(\sin x-x\cos x)'=x\sin x\ge0$.  Taking the limit at zero
proves the claim; continuity handles $\pi/2$.

The following explicit harmonic remainder will also be useful:
\begin{equation}\label{polyretry:harmonic-remainder}
 D_m:=2H_{2m}-H_m-\log m-\gamma-2\log2
 \ge\frac1{6(2m+1)^2}>0\qquad(m\ge1).
\end{equation}
Indeed $D_m\to0$ and
\[
 D_k-D_{k+1}
 =\log\frac{k+1}{k}-\frac2{2k+1}
 \ge\frac{2}{3(2k+1)^3}.
\]
The inequality follows without a logarithm power-series theorem:
for $0\le t<1$,
$\log((1+t)/(1-t))-2t$ has derivative
$2t^2/(1-t^2)\ge2t^2$, so its value is at least $2t^3/3$.
Telescope from $k=m$ to infinity, and bound the decreasing positive
sum below by its integral:
$\int_m^\infty2/[3(2x+1)^3]\,dx=1/[6(2m+1)^2]$.

Let $L=m\pi/N$.  Midpoint concavity on the $m$ equal subintervals
of $[0,L]$, together with
$\int_0^L g(x)\,dx=\log(\sin L/L)$, gives
\begin{align}
 \frac2N\sum_{j=0}^{m-1}\cot\frac{(j+\tfrac12)\pi}{N}
 &\ge\frac2\pi\left(2H_{2m}-H_m+
                  \log\frac{\sin L}{L}\right)\notag\\
 &=B_N+\frac2\pi\bigl(D_m+\log\sin L\bigr).
 \label{polyretry:midpoint-bound}
\end{align}
For even $N$, $L=\pi/2$, so this proves $S_m>B_N$ immediately.

For odd $N\ge3$, put $a=\pi/(2N)$; then $L=\pi/2-a$.
The extra tangent in \eqref{polyretry:central-cot} shows that it is
enough to prove
\[
 D_m\ge E(a):=-\log\cos a-a\tan(a/2).
\]
This small correction has a direct bound.  On $0\le a\le\pi/6$,
\[
 E'(a)=\tfrac12\sec^2(a/2)(\tan a-a),\qquad
 \tan a-a=\int_0^a\tan^2t\,dt\le\frac{a^3}{3\cos^2a}.
\]
Using $\cos^2a,\cos^2(a/2)\ge3/4$ and integrating gives
$E(a)\le2a^4/27<a^4/8$.  Since $\pi^4<192$ and $N\ge3$,
\[
 E(a)<\frac{\pi^4}{128N^4}<\frac1{6N^2}\le D_m.
\]
This completes the odd case as well.  We have proved, for every $N\ge2$,
\begin{equation}\label{polyretry:all-intervals}
 \min_{1\le k<N}S_k>B_N,
 \qquad
 \min_{\text{consecutive root intervals }J}
            \max_{x\in J}\mathcal L_N(x)>B_N.
\end{equation}
For $N=2$ the single sample is $1$; the formula and proof include this
case.  For $N=1$ the Lebesgue constant is $1>B_1$: for example,
$\gamma<1$, $\log(4/\pi)<1/3$, and $\pi>8/3$ suffice.

\subsection{Reduce arbitrary nodes to interval peak comparison}

Call a node set canonical if its first and last nodes are $-1,1$.
For such a set $X$ with $N\ge2$ nodes, write
$M_i(X)=\max_{[x_{i-1},x_i]}\sum_j|\ell_j^X(x)|$.
The comparison theorem proved in the following subsections is
\begin{equation}\label{polyretry:comparison}
 \min_i M_i(Y)\le\max_i M_i(X)
 \quad\text{for any two canonical sets $X,Y$ of the same size.}
\end{equation}
This is a theorem of Bernstein--Erd\H{o}s type, established by
Kilgore and de Boor--Pinkus.  It does not follow merely from the standard
alternation theorem for best polynomial approximation: here the nodes,
and hence the interpolation operator, vary.

Here is the exact reduction, with no loss of constants.  Affinely
stretch the extreme nodes of any arbitrary set $X\subset[-1,1]$ to
$-1,1$.  Its original Lebesgue constant is at least the norm on its
smaller node hull, which equals the norm of the stretched canonical
set.  For $Y$, stretch the Chebyshev roots by
$1/\cos(\pi/(2N))$.  This affine change preserves all the interval
peak values in \eqref{polyretry:all-intervals}.  Therefore
\eqref{polyretry:comparison} and the already proved finite cotangent
estimate yield the catalogue's exact bound $\Lambda_n\ge B_{n+1}$.
For $N=2$ there is only one canonical set, $\{-1,1\}$, so the
comparison is immediate; the case $N=1$ was handled separately.
The following arguments cover $N\ge3$.

For implementation planning, the comparison theorem has a useful
finite algebraic entry point.  Let $n=N-1\ge2$, let
$t_0=-1<t_1<\cdots<t_n=1$, and define the polynomial $F_i$,
$1\le i\le n$, by its nodal values
\[
 F_i(t_j)=
 \begin{cases}(-1)^{i-1-j},&j<i,\\(-1)^{j-i},&j\ge i.\end{cases}
\]
Then $F_i$ equals the Lebesgue function on $(t_{i-1},t_i)$.
Each $F_i$ has a simple real root in every other nodal gap.  Its
degree is $n-1$ or $n$; in the latter case the remaining root is
real, simple, and outside $[-1,1]$.  These statements follow from
the $n-1$ sign changes, the degree bound, and the fact that $F_i>0$
on its own gap.  Rolle's theorem then gives a unique nondegenerate
maximum $\tau_i$ in its own gap.  The positivity on that gap follows
directly from its cardinal representation; since $n\ge2$, at least
one cardinal function there is negative, so the Lebesgue function
is strictly greater than its endpoint value $1$.

Let $\omega(x)=\prod_{j=0}^n(x-t_j)$ and
$q_i(x)=F_i'(x)/(x-\tau_i)$, so $\deg q_i\le n-2$.
Differentiating the nodal equations and using interpolation
uniqueness gives
\[
 \frac{\partial F_i(x)}{\partial t_j}
       =-F_i'(t_j)\ell_j(x),
 \qquad
 \frac{\partial M_i}{\partial t_j}
       =-F_i'(t_j)\ell_j(\tau_i)
       =\frac{\omega(\tau_i)}{\omega'(t_j)}q_i(t_j)
       \quad(1\le j<n).
\]
Thus every maximal minor of the peak Jacobian is, up to explicit
nonzero diagonal factors, an evaluation determinant for $n-1$
polynomials in the $(n-1)$-dimensional space $\mathcal P_{n-2}$.
The following subsections prove the signs and nonsingularity of
these minors by a finite positive-matrix argument, establish
properness of $t\mapsto(M_{i+1}(t)-M_i(t))_{i=1}^{n-1}$ by
explicit cardinal-product estimates, and supply a continuation argument through finite inverse charts.
A single radial path then
proves \eqref{polyretry:comparison}.  Thus each step beyond the
Jacobian identity is supplied rather than assumed.

\subsection{Closing the finite comparison step: a positive Wronskian certificate}
\label{polyretry:finite-comparison}

The remaining finite determinant calculation admits a useful
simplification: a sum of squares replaces differentiation of
interlacing polynomials, and a matrix with positive off-diagonal
entries replaces the sign-variation argument.  We keep
$n\ge2$, $-1=t_0<\cdots<t_n=1$, and the polynomials $F_i$ and
their unique nondegenerate interval maxima $\tau_i$ from above.
Thus $\tau_1<\cdots<\tau_n$,
$F_i(\tau_i)>1$, and $F_i''(\tau_i)<0$.

\paragraph{A finite rational identity.}
Fix $1\le i<j\le n$, and put
\[
 s=(-1)^{j-i},\qquad
 u=\frac{F_i+sF_j}{2},\qquad
 v=\frac{F_i-sF_j}{2}.
\]
The nodal values are particularly simple:
\[
 \begin{array}{c|cc}
 &k<i\text{ or }k\ge j&i\le k<j\\ \hline
 u(t_k)&F_i(t_k)&0\\
 v(t_k)&0&F_i(t_k).
 \end{array}
\]
The polynomial $v$ vanishes at
$t_0,\ldots,t_{i-1},t_j,\ldots,t_n$ and changes sign in each
of the $j-i-1$ gaps between $t_i,\ldots,t_{j-1}$.
These give exactly $n$ distinct real roots.  Since $\deg v\le n$
and $v(t_i)=1$, its degree is $n$ and all these roots are simple.
Call them $\beta_1<\cdots<\beta_n$; in particular $\beta_1=t_0$.

The polynomial $u$ has the $j-i$ roots
$t_i,\ldots,t_{j-1}$ and a sign-changing root in each gap
$(t_{r-1},t_r)$ with $r<i$ or $r>j$.  These are $n-1$
distinct roots, exactly one between every consecutive pair of
the $\beta_r$.  They are simple and there are no further roots
inside $[t_0,t_n]$.  Here are the multiplicity details: in each
outer gap the endpoint values have opposite signs, so the total
root multiplicity there is odd.  In $(t_{i-1},t_j)$ the endpoint
values of $u$ have sign ratio $(-1)^{j-i}$, so the total root
multiplicity there has parity $j-i$.  We already have $n-1$
distinct roots and degree at most $n$; putting an additional root
or multiplicity in any of these regions would violate its parity.
The other nodes have nonzero values.  Thus any remaining root,
if $\deg u=n$, lies outside $[t_0,t_n]$.

Polynomial interpolation, or division followed by interpolation,
now gives the exact finite partial-fraction identity
\begin{equation}\label{polyretry:positive-residues}
 \frac{u(x)}{v(x)}
   =c+\sum_{r=1}^n\frac{a_r}{x-\beta_r},
 \qquad a_r=\frac{u(\beta_r)}{v'(\beta_r)}>0.
\end{equation}
To verify the signs, $u$ has one simple root between successive
$\beta_r$, so the signs of $u(\beta_r)$ alternate, just as those
of $v'(\beta_r)$ do.  Hence all residues have the same sign.
At $t_i$, precisely $i$ roots of $v$ lie to the left; since
$v(t_i)=1$, its leading coefficient has sign $(-1)^{n-i}$.
Therefore $v'(t_0)$ has sign $(-1)^{i-1}$, which is also the
sign of $u(t_0)=F_i(t_0)$.  This proves strict positivity of
every $a_r$.  The constant $c$ is the ratio of leading
coefficients when $\deg u=n$, and zero otherwise.

Put $v_r(x)=v(x)/(x-\beta_r)$, a polynomial.  Differentiating
\eqref{polyretry:positive-residues} and clearing denominators
gives the polynomial identity
\begin{equation}\label{polyretry:wronskian-sos}
 s\bigl(F_i'(x)F_j(x)-F_i(x)F_j'(x)\bigr)
       =2\sum_{r=1}^n a_r v_r(x)^2>0
       \qquad(x\in\R).
\end{equation}
For the identity it suffices to check points where $v\ne0$,
then use polynomial equality.  Strict positivity holds at all
real points: away from the roots every $v_r$ is nonzero; at
$\beta_r$ the term $v_r(\beta_r)=v'(\beta_r)$ is nonzero.
This is a finite Gram certificate for the Wronskian sign.

\paragraph{Signs of the divided derivative polynomials.}
Recall $q_i(x)=F_i'(x)/(x-\tau_i)$, with $\deg q_i\le n-2$.
Evaluation of \eqref{polyretry:wronskian-sos} at $\tau_j$
and at $\tau_i$, using $F_j(\tau_j),F_i(\tau_i)>0$, gives
\[
 \operatorname{sgn}F_i'(\tau_j)=(-1)^{j-i},\qquad
 \operatorname{sgn}F_j'(\tau_i)=-(-1)^{j-i}
       \quad(i<j).
\]
Dividing by the sign of $\tau_j-\tau_i$ yields
\begin{equation}\label{polyretry:q-signs}
 (-1)^{i+j}q_i(\tau_j)>0\quad(i\ne j),
 \qquad q_i(\tau_i)=F_i''(\tau_i)<0.
\end{equation}
In particular, this calculation does not require a separate
theorem saying that derivatives preserve strict interlacing.

\paragraph{A finite matrix proves rank and a positive normal.}
\label{polyretry:peak-laplacian}
Let
\[
 b_r=\frac1{\prod_{k\ne r}|\tau_r-\tau_k|}>0,\qquad
 B_{ri}=(-1)^{r+i}q_i(\tau_r),\qquad A=\diag(b)B.
\]
The off-diagonal entries of $A$ are strictly positive, and its
diagonal entries are strictly negative.  For every polynomial
$q$ of degree at most $n-2$, comparison of leading coefficients
in its Lagrange formula at the $n$ points $\tau_r$ gives
\[
 \sum_{r=1}^n
 \frac{q(\tau_r)}{\prod_{k\ne r}(\tau_r-\tau_k)}=0.
\]
The denominator has sign $(-1)^{n-r}$.  Applying this identity
to each $q_i$ proves that every column of $A$ sums to zero.

Delete the same row and column $k$ from $A$.  Each remaining
column $i$ satisfies strict diagonal dominance:
\[
 |A_{ii}|=\sum_{r\ne i}A_{ri}
          >\sum_{r\ne i,k}A_{ri}.
\]
The resulting square matrix is invertible.  Indeed its
transpose is strictly row diagonally dominant; for a nonzero
kernel vector choose a coordinate of maximum absolute value,
and the triangle inequality contradicts the strict dominance
in that row.  It follows that any $n-1$ columns of $B$ are
linearly independent, and hence any $n-1$ of the polynomials
$q_1,\ldots,q_n$ are linearly independent.

The same matrix gives a positive nullvector explicitly.
Let $A^{(k)}$ denote its $k$th deleted principal submatrix.
Hold its negative diagonal fixed and multiply its off-diagonal
entries by $s\in[0,1]$.  Strict diagonal dominance persists,
so the determinant never vanishes.  At $s=0$ its sign is
$(-1)^{n-1}$.  Therefore
\begin{equation}\label{polyretry:laplacian-cofactor}
 \pi_k:=(-1)^{n-1}\det A^{(k)}>0\qquad(1\le k\le n).
\end{equation}
Since $A$ has rank $n-1$ and zero column sums, its left
nullspace is the span of $(1,\ldots,1)$.  The identities
$A\operatorname{adj}(A)=\operatorname{adj}(A)A=0$ show that
every row of its adjugate is constant.  The constant in row
$i$ is its diagonal entry $\det A^{(i)}$.  Thus every column
of the adjugate is $(-1)^{n-1}\pi$, and $A\pi=0$.
Undoing the row scalings gives
$\sum_i(-1)^i\pi_iq_i(\tau_r)=0$ for every $r$.
This polynomial has degree at most $n-2$ and $n$ distinct
zeros, so
\begin{equation}\label{polyretry:positive-q-relation}
 \sum_{i=1}^n(-1)^i\pi_iq_i=0.
\end{equation}

Evaluation at the $n-1$ distinct interior nodes
$t_1,\ldots,t_{n-1}$ is an isomorphism on
$\mathcal P_{n-2}$.  Consequently every maximal minor of
$(q_i(t_j))_{1\le i\le n,\,1\le j<n}$ is nonzero.
The explicit diagonal factors in the earlier peak-Jacobian
formula then prove the same fact for
\begin{equation}\label{polyretry:peak-jacobian}
 J_{ij}=\frac{\partial M_i}{\partial t_j}
       =\frac{\omega(\tau_i)}{\omega'(t_j)}q_i(t_j).
\end{equation}
In particular $\rank J=n-1$.  Since
$\operatorname{sgn}\omega(\tau_i)=(-1)^{n-i+1}$,
\eqref{polyretry:positive-q-relation} proves directly that
\[
 \widetilde w_i=\frac{\pi_i}{|\omega(\tau_i)|}>0,
 \qquad \widetilde w^TJ=0.
\]
Indeed
$\widetilde w_i\omega(\tau_i)=(-1)^{n+1}(-1)^i\pi_i$,
with a common sign independent of $i$.  Normalizing gives
the explicit positive normal
\begin{equation}\label{polyretry:positive-normal}
 w_i=\frac{\widetilde w_i}{\sum_r\widetilde w_r}>0,
 \qquad\sum_iw_i=1,\qquad w^TJ=0.
\end{equation}
This argument uses no extremization over node sets.

\paragraph{A quantitative boundary estimate.}
The following properness argument uses only products of
distances.  Let $\delta$ be the smallest node gap and let $y$
be the midpoint of a largest gap.  Since the total length is
$2$, every node has distance at least $1/n$ from $y$.
For $x$ in a gap of length $\delta$, every cardinal numerator
$\prod_{r\ne j}(x-t_r)$ includes at least one endpoint of
that small gap.  That factor has absolute value at most
$\delta$, and each of the other $n-1$ factors is at most $2$.
Cancellation of the cardinal denominators gives
\[
 |\ell_j(x)|\le C_n\delta\,|\ell_j(y)|,
 \qquad C_n=2^{n-1}n^n.
\]
If $M_{\mathrm{s}}$ and $M_{\mathrm{l}}$ are the peaks in
the chosen smallest and largest gaps, respectively, summing
and maximizing yields
\[
 1\le M_{\mathrm{s}}\le C_n\delta M_{\mathrm{l}}.
\]
Whenever $C_n\delta<1$, therefore,
\begin{equation}\label{polyretry:properness}
 \max_i M_i\ge\frac1{C_n\delta},\qquad
 \max_i M_i-\min_i M_i\ge\frac1{C_n\delta}-1.
\end{equation}
Thus $\max_i M_i$ tends to infinity at the boundary of the
open node simplex.  So does the norm of
$\Gamma(t)=(M_{i+1}(t)-M_i(t))_{i=1}^{n-1}$, because the
range of the $M_i$ is at most the sum of the absolute adjacent
differences.  In particular these functions have compact
sublevel sets, and $\Gamma$ is proper.


\subsection{Only radial continuation is needed}
\label{polyretry:global-comparison}
\label{polyretry:proper-inverse}
Here is a replacement for the earlier full proper-map theorem. It does
not need surjectivity onto every vector of peak differences, nor a
two-parameter homotopy or covering-space theorem.

\begin{lemma}[Radial continuation and a discrete endpoint map]
\label{lem:radial}
Let $U\subset\R^d$ be nonempty, open and connected, and let
$M:U\to\R^{d+1}$ be $C^1$. Suppose $DM$ has rank $d$ everywhere and
has a strictly positive left nullvector. Set
\[
 \Gamma(u)=(M_{i+1}(u)-M_i(u))_{i=1}^d.
\]
Assume $\{u\in U:\norm{\Gamma(u)}\le R\}$ is compact for every
$R\ge0$. Then there is one equal-peak point $u_*$, and, writing its
common peak as $\lambda_*$,
\[
 \min_i M_i(u)\le\lambda_*\le\max_i M_i(u)\qquad(u\in U).
\]
In particular $\min_iM_i(v)\le\max_iM_i(u)$ for all $u,v\in U$.
\end{lemma}
\begin{proof}
Normalize the positive left nullvector to $w_i>0$, $\sum_iw_i=1$.
Rank $d$ makes it unique and continuous, for example by the nonzero
maximal-minor formulas. If $D\Gamma(u)z=0$, then $DM(u)z=c\mathbf1$.
Multiplication by $w^T$ gives $c=0$ and rank gives $z=0$. Thus
$D\Gamma$ is invertible and the ordinary inverse function theorem gives
local $C^1$ inverse charts.

\emph{Existence of one lifted segment.} Fix $u\in U$, put $h=\Gamma(u)$,
and lift $s\mapsto sh$, starting at $z(1)=u$, toward $s=0$, using a local
inverse. Two lifts with the same starting value agree wherever both
exist: equality is open by local injectivity and closed by continuity.
Consequently the local lifts concatenate uniquely on a maximal interval.
On this interval they lie in the compact set
$K=\{v:\norm{\Gamma(v)}\le\norm h\}$. There is a finite bound
$C=\max_K\norm{D\Gamma(v)^{-1}}$. Differentiating gives
\[
 z'(s)=D\Gamma(z(s))^{-1}h,\qquad\norm{z'(s)}\le C\norm h.
\]
Thus $z$ is Lipschitz and has a limit at any finite left endpoint of its
domain. The limit lies in $K\subset U$, and its inverse chart extends
the lift. A positive left endpoint is therefore impossible, and the lift
exists on $[0,1]$. Its endpoint $R(u)=z(0)$ lies in $\Gamma^{-1}(0)$.

\emph{The endpoint map is locally constant.} Fix this lifted path.
Its compact image is covered by inverse charts. Partition $[0,1]$ into
finitely many closed subintervals, each carried by $z$ into one chart
and by $s\mapsto sh$ into its open image. Such a partition follows by
first taking slightly smaller chart neighborhoods and then a Lebesgue
number for their inverse images in $[0,1]$. At a shared endpoint the two
chosen charts overlap around the lifted point.

For a sufficiently close starting point $u'$, the target segment
$s\Gamma(u')$ is uniformly close to $sh$. In the first chart (starting
from $s=1$) its inverse branch starts at $u'$. Each compact segment stays
in the corresponding chart image for sufficiently small perturbations.
Continuity of the inverse branch at the shared endpoint keeps it in the
overlap with the next chart. Shrink the neighborhood of $u$ successively
for these finitely many transitions. This constructs the entire perturbed
lift using the same finite chart chain. At $s=0$ its target is exactly
zero, so its value is the same inverse-branch value $R(u)$.
Uniqueness of lifting shows $R(u')=R(u)$.

Connectedness of $U$ makes this locally constant map constant. Every
point $v$ with $\Gamma(v)=0$ has the constant radial lift, hence $R(v)=v$.
Since at least one such point was constructed, the fiber consists of
exactly one point $u_*$.

\emph{The scalar identity giving the bound.} For the lift from $u$ put
$a_i=M_i(u)$. Its adjacent differences imply
$M_i(z(s))=c(s)+sa_i$ for a common $C^1$ function $c$, with
$c(1)=0,c(0)=\lambda_*$. Differentiate and multiply by the positive
normal:
\[
 c'(s)=-\sum_iw_i(z(s))a_i,\qquad
 \lambda_*=\int_0^1\sum_iw_i(z(s))a_i\,ds.
\]
The integrand is a convex combination of the original peaks. This proves
both inequalities.
\end{proof}

\paragraph{Application and exact scope.}
Take the open node simplex
$U=\{-1<t_1<\cdots<t_{n-1}<1\}$, with $d=n-1$. The simple critical
points $\tau_i$ vary $C^1$ by the implicit function theorem, so $M$ is
$C^1$. The polynomial and matrix identities above prove its rank and
positive normal. The collision estimate shows that bounded $\Gamma$
keeps the smallest gap bounded away from zero. Such a locus is closed
inside a compact closed simplex with this positive gap bound, and hence
compact. Lemma~\ref{lem:radial} therefore proves
\eqref{polyretry:comparison}. Together with the explicit Chebyshev-root
samples it proves the universal bound, including $n=0,1$ as treated above.

The analytic work left to implement is a local inverse theorem, continuation
of a single segment in a compact set, and a finite chart-chain argument.
The document does not substitute a weaker projection lower bound or an
unspecified $O(1)$ term. A purely finite proof for \emph{arbitrary} grids
would still be a further improvement.


\section{Bounded variation, finite integral identities, and quadrature}
\label{sec:bv}
\subsection{Precise regularity and a step-function identity}
Let $\nu\ge0$ be an integer. The regularity hypothesis is the existence
of functions $f_0=f,f_1,\ldots,f_\nu$ on $[-1,1]$ such that
\begin{equation}\label{eq:regularity}
 f_j\text{ is absolutely continuous and }f_j'=f_{j+1}
 \text{ a.e. for }j<\nu,
 \qquad \Var_{[-1,1]}(f_\nu)\le V<\infty.
\end{equation}
The last function is an actual bounded-variation representative. Values
at exceptional derivative points are irrelevant to integration; $V$ is
the variation of the representative chosen. If $\nu=0$, $f$ may be
discontinuous, and only the coefficient bound below is asserted.
For $\nu\ge1$, $f$ is continuous. No global classical derivative of the
BV function is assumed.

\begin{lemma}[Finite summation by parts for a BV function]\label{lem:bvbase}
If $h$ has variation at most $V$ on $[0,\pi]$ and $k\ge1$ is an integer,
\[
 \left|\int_0^\pi h(\theta)\cos(k\theta)\,d\theta\right|\le V/k.
\]
\end{lemma}
\begin{proof}
For an integer $J\ge1$ put $t_j=j\pi/J$ and let $h_J=h(t_j)$ on
$[t_j,t_{j+1})$. Additivity of variation on adjacent intervals gives
\[
 \int_0^\pi|h-h_J|\le\frac\pi J
       \sum_{j=0}^{J-1}\Var_{[t_j,t_{j+1}]}h\le\frac{\pi V}{J}.
\]
BV functions are bounded and measurable, so these are ordinary integrals.
Writing $s_j=\sin(kt_j)$, a finite identity gives
\begin{equation}\label{eq:sbp}
 \int_0^\pi h_J(\theta)\cos(k\theta)\,d\theta
 =\frac1k\sum_{j=0}^{J-1}h(t_j)(s_{j+1}-s_j)
 =-\frac1k\sum_{j=1}^{J-1}(h(t_j)-h(t_{j-1}))s_j.
\end{equation}
The boundary terms vanish since $s_0=s_J=0$. The last absolute value is
at most $V/k$. The displayed $L^1$ error tends to zero, proving the claim.
\end{proof}
Composition with the decreasing map $\cos:[0,\pi]\to[-1,1]$ does not
increase variation. Thus \eqref{eq:coeff} gives, without Stieltjes integration,
\begin{equation}\label{eq:bvzero}
 |a_k(g)|\le\frac{2\Var(g)}{\pi k}\qquad(k\ge1).
\end{equation}

\subsection{One ordinary integration by parts gives the recursion}
\begin{lemma}\label{lem:recurrence}
If $u$ is absolutely continuous and $u'=v$ a.e. with a bounded measurable
representative $v$, then
\begin{align}
 a_k(u)&=\frac2{\pi k}\int_{-1}^1v(x)\sin(k\arccos x)\,dx,
                  &&k\ge1,\label{eq:finite-ibp}\\
 a_k(u)&=\frac{a_{k-1}(v)-a_{k+1}(v)}{2k},&&k\ge2,\label{eq:recurrence}\\
 a_1(u)&=a_0(v)-\tfrac12a_2(v).&&\label{eq:recurrence-one}
\end{align}
\end{lemma}
\begin{proof}
Put $F_k(x)=\sin(k\arccos x)$. It vanishes at $\pm1$, and in the
interior $F_k'=-kT_k(x)/\sqrt{1-x^2}$. This derivative is integrable.
The fundamental theorem on trimmed intervals, followed by a limit to the
endpoints, gives $F_k(x)=\int_{-1}^x F_k'(t)\,dt$. Thus $F_k$ is
absolutely continuous. Ordinary integration by parts with $u$ gives
\eqref{eq:finite-ibp}, with no endpoint term. Substitute $x=\cos\theta$
and use $2\sin\theta\sin(k\theta)=\cos((k-1)\theta)-\cos((k+1)\theta)$.
The constant coefficient has half the normalization of positive modes,
which explains the separate formula \eqref{eq:recurrence-one}.
\end{proof}

\begin{theorem}[Coefficient decay with the original constant]\label{thm:bv}
Under \eqref{eq:regularity},
\begin{equation}\label{eq:bvdecay}
 |a_k(f)|\le\frac{2V}{\pi\,k(k-1)\cdots(k-\nu)}\qquad(k\ge\nu+1).
\end{equation}
\end{theorem}
\begin{proof}
At $\nu=0$ this is \eqref{eq:bvzero}. For $\nu\ge1$ use
\eqref{eq:recurrence} and apply induction to $f_1$ at $k-1$ and $k+1$.
Both bounds are no larger than
$2V/[\pi(k-1)(k-2)\cdots(k-\nu)]$. Divide their sum by $2k$.
Every derivative representative in this argument is bounded: it is AC,
hence continuous, until the final BV level. The recurrence at $k=1$ is
never used in this induction, since $k\ge\nu+1\ge2$ at a recursive step.
\end{proof}

\subsection{An exact tail telescope and identification of the series}
For $\nu\ge1$ the scalar identity
\[
 \frac1{k(k-1)\cdots(k-\nu)}
 =\frac1\nu\left(\frac1{(k-1)\cdots(k-\nu)}
                 -\frac1{k\cdots(k-\nu+1)}\right)
\]
gives, first for a finite sum and then by its explicit remainder,
\begin{equation}\label{eq:tail}
 \sum_{k=n+1}^\infty |a_k(f)|
 \le\frac{2V}{\pi\nu\,n(n-1)\cdots(n-\nu+1)}\qquad(n\ge\nu).
\end{equation}
In particular the series $g=\sum_{k\ge0}a_k(f)T_k$ converges absolutely
and uniformly. Here is the identification $g=f$, so that the estimate is
not merely a bound on a formal series. Cosine orthogonality shows that
$f-g$ has all Chebyshev coefficients zero. The $T_k$ span the polynomials,
so
\[
 \int_{-1}^1\frac{(f-g)(x)p(x)}{\sqrt{1-x^2}}\,dx=0
 \quad\text{for every polynomial }p.
\]
Approximate the continuous function $f-g$ uniformly by polynomials.
The weight has finite integral $\pi$, and passing to the limit gives a
zero integral of $(f-g)^2$ against this positive weight. Continuity gives
$f=g$ throughout $[-1,1]$.

For $n\ge\nu\ge1$, the tail and the finite aliasing identity at the
Lobatto nodes consequently yield
\[
 \norm{f-S_nf}_\infty\le\frac{2V}{\pi\nu\,n(n-1)\cdots(n-\nu+1)},
 \qquad
 \norm{f-I_nf}_\infty\le\frac{4V}{\pi\nu\,n(n-1)\cdots(n-\nu+1)}.
\]
Indeed $I_nT_k=T_{r_n(k)}$, where $r_n$ is the fold defined below,
and $\norm{T_k-T_{r_n(k)}}_\infty\le2$.

\subsection{Quadrature as a finite functional on coefficients}
Write $I(f)=\int_{-1}^1f(x)\,dx$. Let $C_n=I\circ I_n$ be the
$(n+1)$-point Clenshaw--Curtis rule ($n\ge1$), and let $G_n$ be the
$(n+1)$-point Gauss--Legendre rule ($n\ge0$). The symbol $C_n$ in this
subsection denotes the rule, not the cotangent constant of
Theorem~\ref{thm:lobatto}. All following quadrature bounds are for
Lebesgue measure $dx$ on $[-1,1]$.

\paragraph{The exact moment vector.}
Direct integration gives
\begin{equation}\label{eq:moments}
 m_k:=I(T_k)=
 \begin{cases}2/(1-k^2),&k\text{ even},\\0,&k\text{ odd}.
 \end{cases}
\end{equation}
For $k\ge2$ an antiderivative is
$T_{k+1}/[2(k+1)]-T_{k-1}/[2(k-1)]$; differentiating it in cosine
coordinates verifies the identity. The cases $k=0,1$ are direct.
For every finite coefficient vector $a$, therefore,
\begin{equation}\label{eq:functional}
 (I-Q)\left(\sum_{k=0}^K a_kT_k\right)
    =\sum_{k=0}^K a_k\bigl(m_k-Q(T_k)\bigr).
\end{equation}
For absolutely summable coefficients the same formula follows by uniform
convergence whenever $Q$ is a finite evaluation rule.

\paragraph{Gauss exactness and positivity need only a finite Gram matrix.}
Put $N=n+1$. The inner product $\langle p,q\rangle=\int_{-1}^1pq$
is positive definite on $\PP_{N-1}$. Project $X^N$ onto this space and
subtract the projection to obtain the unique monic degree-$N$ polynomial
$q_N$ orthogonal to $\PP_{N-1}$. It has $N$ distinct roots in $(-1,1)$:
otherwise multiply it by the product of its interior sign-changing
roots. This product has degree less than $N$, while its product with
$q_N$ has one constant nonzero sign except at finitely many zeros,
contradicting orthogonality. Degree $N$ then forces exactly $N$ simple
interior roots.

At these nodes take weights $w_j=I(\ell_j)$. Division of a polynomial
of degree at most $2N-1$ by $q_N$ writes it as $q_Ns+r$ with
$s,r\in\PP_{N-1}$. Both integration and the rule annihilate $q_Ns$ and
agree on $r$. Thus the rule is exact through degree $2N-1=2n+1$.
In particular $w_j=G_n(\ell_j^2)=I(\ell_j^2)>0$. Reflection and uniqueness
give $q_N(-x)=(-1)^Nq_N(x)$ and paired equal weights. The weights sum
to two, and the rule annihilates odd polynomials.

\paragraph{Clenshaw--Curtis positivity is a finite telescope.}
Set $\theta_j=j\pi/n$ and let $\delta_j=1/2$ at $j=0,n$ and $1$
otherwise. Finite cosine orthogonality gives
\[
 \ell_j(x)=\frac{2\delta_j}{n}
 \left(\frac12+\sum_{k=1}^{n-1}\cos(k\theta_j)T_k(x)
                     +\frac12\cos(n\theta_j)T_n(x)\right).
\]
This identity can be checked at the $n+1$ nodes by a finite geometric
sum, then by polynomial interpolation uniqueness. Integrating gives,
for $n\ge2$,
\begin{equation}\label{eq:ccweights}
 w_j=\frac{2\delta_j}{n}\left[
 1-\sum_{k=1}^{\lfloor(n-1)/2\rfloor}\frac{2\cos(2k\theta_j)}{4k^2-1}
 -\boldsymbol1_{2\mid n}\frac{\cos(n\theta_j)}{n^2-1}\right].
\end{equation}
Because $2/(4k^2-1)=1/(2k-1)-1/(2k+1)$, the bracket is at least $1/n$
for odd $n$, and at least $n/(n^2-1)$ for even $n$. Thus all weights
are positive. For $n=1$ they are $1,1$. The rule is symmetric, exact
through degree $n$, and its weights sum to two.

\paragraph{Fold the frequencies, not the functions.}
For $n\ge1$ let $s$ be the remainder of $k$ modulo $2n$ and set
\[
 r_n(k)=\min(s,2n-s)\in\{0,\ldots,n\}.
\]
Then $T_k(\cos\theta_j)=T_{r_n(k)}(\cos\theta_j)$, whence
\begin{equation}\label{eq:fold}
 C_n(T_k)=m_{r_n(k)}.
\end{equation}
The fold preserves parity. For even $k>n$, $n\ge2$, one has $k\ge4$,
so $|m_k|\le2/15$ and $|m_{r_n(k)}|\le2$. For even $k>2n+1$,
$n\ge1$, positivity similarly gives $|G_n(T_k)|\le2$. Therefore
\begin{equation}\label{eq:master}
 \begin{aligned}
 |(I-C_n)f|&\le\frac{32}{15}\sum_{\substack{k>n\\k\ {
m even}}}|a_k|,
                      &&n\ge2,\\
 |(I-G_n)f|&\le\frac{32}{15}\sum_{\substack{k>2n+1\\k\ {
m even}}}|a_k|,
                      &&n\ge1.
 \end{aligned}
\end{equation}
These bounds follow from \eqref{eq:functional} and exactness; every odd
mode contributes zero. The restrictions on small $n$ ensure $k\ge4$ and
are part of the statements.

\subsection{Sharp stated constants, with the corrected denominator}
\begin{theorem}[Variation bounds]\label{thm:quad-bv}
Under \eqref{eq:regularity} with $\nu\ge1$,
\begin{align}
 |(I-C_n)f|&\le\frac{32V}{15\pi\nu(n-\nu)^\nu},
                  &&n>\nu,\ n\ge2,\label{eq:cc-bv}\\
 |(I-G_n)f|&\le\frac{32V}{15\pi\nu(2n+1-\nu)^\nu},
                  &&2n+1>\nu,\ n\ge1.\label{eq:g-bv}
\end{align}
\end{theorem}
\begin{proof}
Equation~\eqref{eq:bvdecay} bounds the coefficients by
$2V/[\pi(k-\nu)^{\nu+1}]$. Let $m$ be the first even integer strictly
greater than a cutoff $s$. For $s>\nu$, convexity of
$g(t)=(t-\nu)^{-\nu-1}$ and integration over the length-two intervals
centered at $m,m+2,\ldots$ give
\[
 \sum_{j=0}^J g(m+2j)
 \le\frac12\int_{m-1}^{m+2J+1}g(t)\,dt
 \le\frac1{2\nu(m-1-\nu)^\nu}.
\]
Take $J\to\infty$ and use $m-1\ge s$. Substitute the two cutoffs
$s=n$ and $s=2n+1$ into \eqref{eq:master}.
\end{proof}

\begin{remark}[The denominator is an audited correction]
The denominator in the Gauss bound and the eventual Clenshaw--Curtis
bound is $(2n+1-\nu)^\nu$. The expression
$(n-2\nu-1)^{2\nu+1}$ in the 2013 printing is erroneous, as recorded
in Trefethen's \href{https://people.maths.ox.ac.uk/trefethen/ATAP/}{official
errata} for equations (19.10) and (19.12). Here the corrected denominator
also follows directly from the even-mode cutoff in the proof.
\end{remark}

\begin{corollary}[Geometric coefficient input]
If $f=\sum a_kT_k$ and $|a_k|\le2M\rho^{-k}$ for $k\ge1$, with $\rho>1$,
then
\[
 |(I-C_n)f|\le\frac{64M\rho^{1-n}}{15(\rho^2-1)}\quad(n\ge2),
 \qquad
 |(I-G_n)f|\le\frac{64M\rho^{-2n}}{15(\rho^2-1)}\quad(n\ge1).
\]
If $n$ is even, the first numerator can be replaced by $64M\rho^{-n}$.
\end{corollary}
\begin{proof}
Sum the even geometric tails in \eqref{eq:master}. If the first even
index is $m$, their sum is at most $2M\rho^{2-m}/(\rho^2-1)$.
\end{proof}
This corollary isolates exactly the analytic input: a function bounded
by $M$ and holomorphic in the open Bernstein ellipse $E_\rho$ has the
stated coefficient bound. One can reuse NLAlib's open-ellipse coefficient
and series theorems. Alternatively apply the Laurent coefficient estimate
to $f((z+z^{-1})/2)$ on circles $|z|=r<\rho$ and let $r\uparrow\rho$;
the Laurent symmetry makes $a_k$ twice the positive Laurent coefficient.

\subsection{An explicit eventual threshold from finite alias blocks}
The point of the next argument is that the eventual Clenshaw--Curtis
bound needs no separate asymptotic quadrature theorem.

\begin{lemma}[Block estimate for a coefficient envelope]\label{lem:block}
If $\nu\ge1$, $n\ge2$, and $|a_k|\le A k^{-\nu-1}$ for all $k>n$, then
\begin{equation}\label{eq:blockrate}
 |(I-C_n)f|\le 7A n^{-\nu-1}.
\end{equation}
\end{lemma}
\begin{proof}
In any block $2\ell n<k\le2(\ell+1)n$ restricted to even $k$, the
folded residue zero appears once and each positive even residue at most
twice. The same upper bound holds for the initial half-block $n<k\le2n$.
Thus the finite telescope in \eqref{eq:moments} gives
\[
 \sum_{k\text{ in one block, even}}|m_{r_n(k)}|
 \le2+2\sum_{j=1}^\infty\frac2{4j^2-1}=4.
\]
Equivalently first use the finite sum $1-1/(2J+1)\le1$ and then let
$J\to\infty$. Since $\sum_{\ell\ge1}\ell^{-\nu-1}\le1+1/\nu\le2$,
\[
 \sum_{\substack{k>n\\k\ {
m even}}}k^{-\nu-1}|m_{r_n(k)}|
 \le\frac4{n^{\nu+1}}+
       \frac4{(2n)^{\nu+1}}\sum_{\ell\ge1}\ell^{-\nu-1}
 \le\frac6{n^{\nu+1}}.
\]
For the unfurled moments, $|m_k|\le(8/3)k^{-2}$ for $k\ge2$, so
\[
 \sum_{k>n}k^{-\nu-1}|m_k|
 \le\frac8{3(\nu+2)n^{\nu+2}}\le\frac1{n^{\nu+1}}.
\]
Combine these estimates with \eqref{eq:fold} and \eqref{eq:functional}.
\end{proof}

\begin{lemma}[A finite product replaces an exponential loss]\label{lem:product}
Let $s_\nu=\nu(\nu+1)/2$. If $n\ge\nu$ and $n+1>s_\nu$, then the BV
coefficient envelope for every $k>n$ is
\begin{equation}\label{eq:envelope}
 |a_k(f)|\le
 \frac{2V}{\pi(1-s_\nu/(n+1))}\,k^{-\nu-1}.
\end{equation}
\end{lemma}
\begin{proof}
For $0\le u_j\le1$, finite induction gives
$\prod_j(1-u_j)\ge1-\sum_j u_j$. Apply this with $u_j=j/k$,
$0\le j\le\nu$. Then
\[
 k(k-1)\cdots(k-\nu)
 =k^{\nu+1}\prod_{j=0}^\nu(1-j/k)
 \ge k^{\nu+1}(1-s_\nu/k)
 \ge k^{\nu+1}(1-s_\nu/(n+1))>0.
\]
Use \eqref{eq:bvdecay}. This avoids replacing every factor by $k/2$.
\end{proof}

\begin{theorem}[An explicit eventual bound, with a smaller threshold]
\label{thm:threshold}
Under \eqref{eq:regularity}, for $\nu\ge1$ put
\begin{equation}\label{eq:threshold}
 n_0(\nu)=\frac{\nu(\nu+1)}2+
            \left\lceil\frac{105}{16}\nu2^\nu\right\rceil.
\end{equation}
For every integer $n\ge n_0(\nu)$,
\begin{equation}\label{eq:eventual}
 |(I-C_n)f|\le\frac{32V}{15\pi\nu(2n+1-\nu)^\nu}.
\end{equation}
More explicitly, whenever $n\ge\max(2,\nu)$ and $n+1>s_\nu$,
\begin{equation}\label{eq:strongrate}
 |(I-C_n)f|\le
 \frac{14V}{\pi(1-s_\nu/(n+1))n^{\nu+1}}.
\end{equation}
\end{theorem}
\begin{proof}
Combine Lemmas~\ref{lem:block} and \ref{lem:product} to obtain
\eqref{eq:strongrate}. The threshold \eqref{eq:threshold} implies all
its hypotheses and
\[
 n(1-s_\nu/(n+1))=n-s_\nu+\frac{s_\nu}{n+1}
 \ge\frac{105}{16}\nu2^\nu.
\]
This is exactly the inequality needed to bound the right side of
\eqref{eq:strongrate} by $32V/[15\pi\nu(2n)^\nu]$.
Since $0<2n+1-\nu\le2n$, this is at most the desired right side of
\eqref{eq:eventual}. No division by $V$ is needed, so $V=0$ is included.
\end{proof}
For illustration, $n_0(1)=15$, $n_0(2)=56$, $n_0(3)=164$, and
$n_0(4)=430$. These are sufficient thresholds, not assertions of optimal
thresholds. The constants $32/15$ and the corrected denominator of the
target bound are unchanged. The stronger-rate estimate
\eqref{eq:strongrate} additionally displays an explicit $n^{-\nu-1}$
bound.

\subsection{Interfaces for a later formalization}
The natural handoff consists of four short dependency chains:
\begin{enumerate}
\item Fixed Vandermonde inverse $\to$ coefficient compactness $\to$ local
root multiplicity and cancellation; separately, rank--nullity and the
signed Lagrange annihilator $\to$ rational alternation and uniqueness.
\item Finite trigonometric cardinal identities $\to$ the cotangent
pair inequality $\to$ the $2/\pi$ bound; finite cosine convolution and
period-by-period integration $\to$ the $4/\pi^2$ bound.
\item For the universal lower bound: polynomial partial fractions and
a positive matrix $\to$ rank and a positive normal; the finite collision
estimate $\to$ compact continuation. The last step remains global.
\item Finite summation by parts and one AC integration by parts $\to$
the coefficient recurrence; exact moment/fold vectors $\to$ the quadrature
bounds and the explicit threshold.
\end{enumerate}
Existing interfaces worth reusing are Lagrange interpolation and
leading-coefficient sums, polynomial gcd and root counting, matrix
adjugates and strict diagonal dominance, variation on adjacent intervals,
and the existing Chebyshev alias and analytic coefficient interfaces.
The requested mathematical statements are proved here, but assembling
these interfaces still requires Lean work; a finite-looking proof is
not by itself evidence that the formalization will be short.

\begin{thebibliography}{9}
\bibitem{dp} C. de Boor and A. Pinkus,
\emph{Proof of the conjectures of Bernstein and Erd\H{o}s concerning the
optimal nodes for polynomial interpolation}, J. Approx. Theory 24 (1978),
289--303. \url{https://doi.org/10.1016/0021-9045(78)90014-X}.
The peak-comparison architecture motivates Section~\ref{sec:universal};
the finite certificates and the radial-continuation lemma are supplied
explicitly in this document.
\bibitem{brutman} L. Brutman, \emph{On the Lebesgue Function for Polynomial
Interpolation}, SIAM J. Numer. Anal. 15 (1978), 694--704.
\url{https://doi.org/10.1137/0715046}.
\bibitem{ez} H. Ehlich and K. Zeller, \emph{Auswertung der Normen von
Interpolationsoperatoren}, Math. Ann. 164 (1966), 105--112.
\url{https://doi.org/10.1007/BF01429047}.
\bibitem{atap} L. N. Trefethen, \emph{Approximation Theory and Approximation
Practice}, SIAM; author's errata for the 2013 edition, equations (19.10)
and (19.12): \url{https://people.maths.ox.ac.uk/trefethen/ATAP/}.
\end{thebibliography}

\end{document}
